\documentclass[10pt,a4paper]{amsart}
\usepackage{amsmath,amssymb,amsthm,graphicx}
% For commutative diagrams we use Xy-pic
\usepackage[all]{xy}

% We use 2cell for 2-commutative diagrams.
\xyoption{2cell}
% Theorem environments.
%
\theoremstyle{plain}
\newtheorem{theorem}[subsection]{Theorem}
\newtheorem{proposition}[subsection]{Proposition}
\newtheorem{lemma}[subsection]{Lemma}

\theoremstyle{definition}
\newtheorem{definition}[subsection]{Definition}
\newtheorem{example}[subsection]{Example}
\newtheorem{exercise}[subsection]{Exercise}
\newtheorem{situation}[subsection]{Situation}

\theoremstyle{remark}
\newtheorem{remark}[subsection]{Remark}
\newtheorem{remarks}[subsection]{Remarks}

\numberwithin{equation}{subsection}

% Macros
%
\def\lim{\mathop{\mathrm{lim}}\nolimits}
\def\colim{\mathop{\mathrm{colim}}\nolimits}
\def\Spec{\mathop{\mathrm{Spec}}}
\def\Hom{\mathop{\mathrm{Hom}}\nolimits}
\def\Ext{\mathop{\mathrm{Ext}}\nolimits}
\def\SheafHom{\mathop{\mathcal{H}\!\mathit{om}}\nolimits}
\def\SheafExt{\mathop{\mathcal{E}\!\mathit{xt}}\nolimits}
\def\Sch{\mathit{Sch}}
\def\Mor{\mathop{\mathrm{Mor}}\nolimits}
\def\Ob{\mathop{\mathrm{Ob}}\nolimits}
\def\Sh{\mathop{\mathit{Sh}}\nolimits}
\def\NL{\mathop{N\!L}\nolimits}
\def\CH{\mathop{\mathrm{CH}}\nolimits}
\def\proetale{{pro\text{-}\acute{e}tale}}
\def\etale{{\acute{e}tale}}
\def\QCoh{\mathit{QCoh}}
\def\Ker{\mathop{\mathrm{Ker}}}
\def\Im{\mathop{\mathrm{Im}}}
\def\Coker{\mathop{\mathrm{Coker}}}
\def\Coim{\mathop{\mathrm{Coim}}}

% Boxtimes
%
\DeclareMathSymbol{\boxtimes}{\mathbin}{AMSa}{"02}

%
% Macros for moduli stacks/spaces
%
\def\QCohstack{\mathcal{QC}\!\mathit{oh}}
\def\Cohstack{\mathcal{C}\!\mathit{oh}}
\def\Spacesstack{\mathcal{S}\!\mathit{paces}}
\def\Quotfunctor{\mathrm{Quot}}
\def\Hilbfunctor{\mathrm{Hilb}}
\def\Curvesstack{\mathcal{C}\!\mathit{urves}}
\def\Polarizedstack{\mathcal{P}\!\mathit{olarized}}
\def\Complexesstack{\mathcal{C}\!\mathit{omplexes}}
% \Pic is the operator that assigns to X its picard group, usage \Pic(X)
% \Picardstack_{X/B} denotes the Picard stack of X over B
% \Picardfunctor_{X/B} denotes the Picard functor of X over B
\def\Pic{\mathop{\mathrm{Pic}}\nolimits}
\def\Picardstack{\mathcal{P}\!\mathit{ic}}
\def\Picardfunctor{\mathrm{Pic}}
\def\Deformationcategory{\mathcal{D}\!\mathit{ef}}
\newenvironment{reference}{}{}
\newenvironment{slogan}{}{}
\newenvironment{history}{}{}
\usepackage{cleveref}
\crefname{originallemma}{original lemma}{original lemmas}
\Crefname{originallemma}{Original lemma}{Original lemmas}
\makeatletter
\newlabel{varieties-lemma-ample-after-field-extension}{{0BDC}{source}{}{sourcepointer-1}{}}
\newlabel{varieties-lemma-ample-after-field-extension@cref}{{[originallemma][0][]0BDC}{[1][0][]source}{}{}{}}
\makeatother
\crefname{originalproposition}{original proposition}{original propositions}
\Crefname{originalproposition}{Original proposition}{Original propositions}
\makeatletter
\newlabel{proposition-connected-component}{{0B7R}{source}{}{sourcepointer-2}{}}
\newlabel{proposition-connected-component@cref}{{[originalproposition][0][]0B7R}{[1][0][]source}{}{}{}}
\makeatother
\crefname{originallemma}{original lemma}{original lemmas}
\Crefname{originallemma}{Original lemma}{Original lemmas}
\makeatletter
\newlabel{lemma-group-scheme-characteristic-zero-smooth}{{047N}{source}{}{sourcepointer-3}{}}
\newlabel{lemma-group-scheme-characteristic-zero-smooth@cref}{{[originallemma][0][]047N}{[1][0][]source}{}{}{}}
\makeatother
\crefname{originalproposition}{original proposition}{original propositions}
\Crefname{originalproposition}{Original proposition}{Original propositions}
\makeatletter
\newlabel{divisors-proposition-push-down-ample}{{0BD4}{source}{}{sourcepointer-4}{}}
\newlabel{divisors-proposition-push-down-ample@cref}{{[originalproposition][0][]0BD4}{[1][0][]source}{}{}{}}
\makeatother
\crefname{originallemma}{original lemma}{original lemmas}
\Crefname{originallemma}{Original lemma}{Original lemmas}
\makeatletter
\newlabel{more-morphisms-lemma-quasi-projective}{{0B42}{source}{}{sourcepointer-5}{}}
\newlabel{more-morphisms-lemma-quasi-projective@cref}{{[originallemma][0][]0B42}{[1][0][]source}{}{}{}}
\makeatother
\crefname{originallemma}{original lemma}{original lemmas}
\Crefname{originallemma}{Original lemma}{Original lemmas}
\makeatletter
\newlabel{lemma-reduced-subgroup-scheme-perfect}{{047R}{source}{}{sourcepointer-6}{}}
\newlabel{lemma-reduced-subgroup-scheme-perfect@cref}{{[originallemma][0][]047R}{[1][0][]source}{}{}{}}
\makeatother
\crefname{originallemma}{original lemma}{original lemmas}
\Crefname{originallemma}{Original lemma}{Original lemmas}
\makeatletter
\newlabel{lemma-reduced-group-scheme-prefect-field-characteristic-p-smooth}{{047P}{source}{}{sourcepointer-7}{}}
\newlabel{lemma-reduced-group-scheme-prefect-field-characteristic-p-smooth@cref}{{[originallemma][0][]047P}{[1][0][]source}{}{}{}}
\makeatother
\crefname{originallemma}{original lemma}{original lemmas}
\Crefname{originallemma}{Original lemma}{Original lemmas}
\makeatletter
\newlabel{varieties-lemma-smooth-regular}{{056S}{source}{}{sourcepointer-8}{}}
\newlabel{varieties-lemma-smooth-regular@cref}{{[originallemma][0][]056S}{[1][0][]source}{}{}{}}
\makeatother
\crefname{originallemma}{original lemma}{original lemmas}
\Crefname{originallemma}{Original lemma}{Original lemmas}
\makeatletter
\newlabel{more-algebra-lemma-regular-local-UFD}{{0AG0}{source}{}{sourcepointer-9}{}}
\newlabel{more-algebra-lemma-regular-local-UFD@cref}{{[originallemma][0][]0AG0}{[1][0][]source}{}{}{}}
\makeatother
\crefname{originallemma}{original lemma}{original lemmas}
\Crefname{originallemma}{Original lemma}{Original lemmas}
\makeatletter
\newlabel{divisors-lemma-complement-open-affine-effective-cartier-divisor}{{0BCW}{source}{}{sourcepointer-10}{}}
\newlabel{divisors-lemma-complement-open-affine-effective-cartier-divisor@cref}{{[originallemma][0][]0BCW}{[1][0][]source}{}{}{}}
\makeatother
\crefname{originallemma}{original lemma}{original lemmas}
\Crefname{originallemma}{Original lemma}{Original lemmas}
\makeatletter
\newlabel{lemma-group-scheme-over-field-separated}{{047L}{source}{}{sourcepointer-11}{}}
\newlabel{lemma-group-scheme-over-field-separated@cref}{{[originallemma][0][]047L}{[1][0][]source}{}{}{}}
\makeatother
\crefname{originallemma}{original lemma}{original lemmas}
\Crefname{originallemma}{Original lemma}{Original lemmas}
\makeatletter
\newlabel{morphisms-lemma-smooth-etale-over-affine-space}{{054L}{source}{}{sourcepointer-12}{}}
\newlabel{morphisms-lemma-smooth-etale-over-affine-space@cref}{{[originallemma][0][]054L}{[1][0][]source}{}{}{}}
\makeatother
\crefname{originalsection}{original section}{original sections}
\Crefname{originalsection}{Original section}{Original sections}
\makeatletter
\newlabel{divisors-section-norms}{{0BCX}{source}{}{sourcepointer-13}{}}
\newlabel{divisors-section-norms@cref}{{[originalsection][0][]0BCX}{[1][0][]source}{}{}{}}
\makeatother
\crefname{originallemma}{original lemma}{original lemmas}
\Crefname{originallemma}{Original lemma}{Original lemmas}
\makeatletter
\newlabel{divisors-lemma-weil-divisor-is-cartier-UFD}{{0AGA}{source}{}{sourcepointer-14}{}}
\newlabel{divisors-lemma-weil-divisor-is-cartier-UFD@cref}{{[originallemma][0][]0AGA}{[1][0][]source}{}{}{}}
\makeatother
\crefname{originallemma}{original lemma}{original lemmas}
\Crefname{originallemma}{Original lemma}{Original lemmas}
\makeatletter
\newlabel{varieties-lemma-rational-equivalence-for-Pic}{{0BEH}{source}{}{sourcepointer-15}{}}
\newlabel{varieties-lemma-rational-equivalence-for-Pic@cref}{{[originallemma][0][]0BEH}{[1][0][]source}{}{}{}}
\makeatother
\crefname{originaldefinition}{original definition}{original definitions}
\Crefname{originaldefinition}{Original definition}{Original definitions}
\makeatletter
\newlabel{divisors-definition-invertible-sheaf-effective-Cartier-divisor}{{01WX}{source}{}{sourcepointer-16}{}}
\newlabel{divisors-definition-invertible-sheaf-effective-Cartier-divisor@cref}{{[originaldefinition][0][]01WX}{[1][0][]source}{}{}{}}
\makeatother
\crefname{originaldefinition}{original definition}{original definitions}
\Crefname{originaldefinition}{Original definition}{Original definitions}
\makeatletter
\newlabel{properties-definition-ample}{{01PS}{source}{}{sourcepointer-17}{}}
\newlabel{properties-definition-ample@cref}{{[originaldefinition][0][]01PS}{[1][0][]source}{}{}{}}
\makeatother
\title{Every algebraic group scheme over a field is quasi-projective}
\author{The Stacks Project Authors}
\date{stacks-project:groupoids.tex:local-snapshot-20261011; GNU FDL 1.2 or later}
\begin{document}
\maketitle
\noindent\textit{Editorial note.} No linearity or affineness assumption on the algebraic group is introduced; the target is quasi-projectivity for arbitrary algebraic group schemes over the given field.\par\medskip
\noindent\textit{Allowed standard prerequisites (editorial).}\begin{itemize}
\item Explicit prerequisite boundary. Algebraic group schemes, connected components, Cartier smoothness in characteristic zero and smoothness of reduced groups over a perfect field, descent of ampleness, regular-local factoriality, etale charts, divisor norms and rational equivalence of Picard classes over an open affine-space base.
\end{itemize}\medskip
\ifdefined\hypertarget\hypertarget{sourcepointer-1}{}\fi
\noindent\textit{Source reference.} Original lemma 0BDC: focused-seeds/AGP/sources/stacks-snapshot-20261011/tags.txt, source lines 14310--14310. Stable Stacks tag for an expressly allowed foundational or notational input; the counted construction remains the two concrete core obligations. The retained local core is not replaced by this pointer. See standard\_prerequisites.\par\medskip
\ifdefined\hypertarget\hypertarget{sourcepointer-2}{}\fi
\noindent\textit{Source reference.} Original proposition 0B7R: focused-seeds/AGP/sources/stacks-snapshot-20261011/tags.txt, source lines 14114--14114. Stable Stacks tag for an expressly allowed foundational or notational input; the counted construction remains the two concrete core obligations. The retained local core is not replaced by this pointer. See standard\_prerequisites.\par\medskip
\ifdefined\hypertarget\hypertarget{sourcepointer-3}{}\fi
\noindent\textit{Source reference.} Original lemma 047N: focused-seeds/AGP/sources/stacks-snapshot-20261011/tags.txt, source lines 5514--5514. Stable Stacks tag for an expressly allowed foundational or notational input; the counted construction remains the two concrete core obligations. The retained local core is not replaced by this pointer. See standard\_prerequisites.\par\medskip
\ifdefined\hypertarget\hypertarget{sourcepointer-4}{}\fi
\noindent\textit{Source reference.} Original proposition 0BD4: focused-seeds/AGP/sources/stacks-snapshot-20261011/tags.txt, source lines 14302--14302. Stable Stacks tag for an expressly allowed foundational or notational input; the counted construction remains the two concrete core obligations. The retained local core is not replaced by this pointer. See standard\_prerequisites.\par\medskip
\ifdefined\hypertarget\hypertarget{sourcepointer-5}{}\fi
\noindent\textit{Source reference.} Original lemma 0B42: focused-seeds/AGP/sources/stacks-snapshot-20261011/tags.txt, source lines 13985--13985. Stable Stacks tag for an expressly allowed foundational or notational input; the counted construction remains the two concrete core obligations. The retained local core is not replaced by this pointer. See standard\_prerequisites.\par\medskip
\ifdefined\hypertarget\hypertarget{sourcepointer-6}{}\fi
\noindent\textit{Source reference.} Original lemma 047R: focused-seeds/AGP/sources/stacks-snapshot-20261011/tags.txt, source lines 5518--5518. Stable Stacks tag for an expressly allowed foundational or notational input; the counted construction remains the two concrete core obligations. The retained local core is not replaced by this pointer. See standard\_prerequisites.\par\medskip
\ifdefined\hypertarget\hypertarget{sourcepointer-7}{}\fi
\noindent\textit{Source reference.} Original lemma 047P: focused-seeds/AGP/sources/stacks-snapshot-20261011/tags.txt, source lines 5516--5516. Stable Stacks tag for an expressly allowed foundational or notational input; the counted construction remains the two concrete core obligations. The retained local core is not replaced by this pointer. See standard\_prerequisites.\par\medskip
\ifdefined\hypertarget\hypertarget{sourcepointer-8}{}\fi
\noindent\textit{Source reference.} Original lemma 056S: focused-seeds/AGP/sources/stacks-snapshot-20261011/tags.txt, source lines 6720--6720. Stable Stacks tag for an expressly allowed foundational or notational input; the counted construction remains the two concrete core obligations. The retained local core is not replaced by this pointer. See standard\_prerequisites.\par\medskip
\ifdefined\hypertarget\hypertarget{sourcepointer-9}{}\fi
\noindent\textit{Source reference.} Original lemma 0AG0: focused-seeds/AGP/sources/stacks-snapshot-20261011/tags.txt, source lines 13178--13178. Stable Stacks tag for an expressly allowed foundational or notational input; the counted construction remains the two concrete core obligations. The retained local core is not replaced by this pointer. See standard\_prerequisites.\par\medskip
\ifdefined\hypertarget\hypertarget{sourcepointer-10}{}\fi
\noindent\textit{Source reference.} Original lemma 0BCW: focused-seeds/AGP/sources/stacks-snapshot-20261011/tags.txt, source lines 14294--14294. Stable Stacks tag for an expressly allowed foundational or notational input; the counted construction remains the two concrete core obligations. The retained local core is not replaced by this pointer. See standard\_prerequisites.\par\medskip
\ifdefined\hypertarget\hypertarget{sourcepointer-11}{}\fi
\noindent\textit{Source reference.} Original lemma 047L: focused-seeds/AGP/sources/stacks-snapshot-20261011/tags.txt, source lines 5512--5512. Stable Stacks tag for an expressly allowed foundational or notational input; the counted construction remains the two concrete core obligations. The retained local core is not replaced by this pointer. See standard\_prerequisites.\par\medskip
\ifdefined\hypertarget\hypertarget{sourcepointer-12}{}\fi
\noindent\textit{Source reference.} Original lemma 054L: focused-seeds/AGP/sources/stacks-snapshot-20261011/tags.txt, source lines 6644--6644. Stable Stacks tag for an expressly allowed foundational or notational input; the counted construction remains the two concrete core obligations. The retained local core is not replaced by this pointer. See standard\_prerequisites.\par\medskip
\ifdefined\hypertarget\hypertarget{sourcepointer-13}{}\fi
\noindent\textit{Source reference.} Original section 0BCX: focused-seeds/AGP/sources/stacks-snapshot-20261011/tags.txt, source lines 14295--14295. Stable Stacks tag for an expressly allowed foundational or notational input; the counted construction remains the two concrete core obligations. The retained local core is not replaced by this pointer. See standard\_prerequisites.\par\medskip
\ifdefined\hypertarget\hypertarget{sourcepointer-14}{}\fi
\noindent\textit{Source reference.} Original lemma 0AGA: focused-seeds/AGP/sources/stacks-snapshot-20261011/tags.txt, source lines 13188--13188. Stable Stacks tag for an expressly allowed foundational or notational input; the counted construction remains the two concrete core obligations. The retained local core is not replaced by this pointer. See standard\_prerequisites.\par\medskip
\ifdefined\hypertarget\hypertarget{sourcepointer-15}{}\fi
\noindent\textit{Source reference.} Original lemma 0BEH: focused-seeds/AGP/sources/stacks-snapshot-20261011/tags.txt, source lines 14350--14350. Stable Stacks tag for an expressly allowed foundational or notational input; the counted construction remains the two concrete core obligations. The retained local core is not replaced by this pointer. See standard\_prerequisites.\par\medskip
\ifdefined\hypertarget\hypertarget{sourcepointer-16}{}\fi
\noindent\textit{Source reference.} Original definition 01WX: focused-seeds/AGP/sources/stacks-snapshot-20261011/tags.txt, source lines 2498--2498. Stable Stacks tag for an expressly allowed foundational or notational input; the counted construction remains the two concrete core obligations. The retained local core is not replaced by this pointer. See standard\_prerequisites.\par\medskip
\ifdefined\hypertarget\hypertarget{sourcepointer-17}{}\fi
\noindent\textit{Source reference.} Original definition 01PS: focused-seeds/AGP/sources/stacks-snapshot-20261011/tags.txt, source lines 2241--2241. Stable Stacks tag for an expressly allowed foundational or notational input; the counted construction remains the two concrete core obligations. The retained local core is not replaced by this pointer. See standard\_prerequisites.\par\medskip
\begin{lemma}
\label{lemma-algebraic-quasi-projective}
Let $k$ be a field. Let $G$ be an algebraic group scheme over $k$.
Then $G$ is quasi-projective over $k$.
\end{lemma}

\begin{proof}
By Varieties, Lemma \ref{varieties-lemma-ample-after-field-extension}
we may assume that $k$ is algebraically closed. Let $G^0 \subset G$
be the connected component of $G$ as in
Proposition \ref{proposition-connected-component}.
Then every other connected component of $G$ has a $k$-rational
point and hence is isomorphic to $G^0$ as a scheme.
Since $G$ is quasi-compact and Noetherian, there are finitely many of these
connected components. Thus we reduce to the case discussed in
the next paragraph.

\medskip\noindent
Let $G$ be a connected algebraic group scheme over an algebraically closed
field $k$. If the characteristic of $k$ is zero, then $G$ is smooth over
$k$ by Lemma \ref{lemma-group-scheme-characteristic-zero-smooth}.
If the characteristic of $k$ is $p > 0$, then we let $H = G_{red}$
be the reduction of $G$. By
Divisors, Proposition \ref{divisors-proposition-push-down-ample}
it suffices to show that $H$ has an ample invertible sheaf.
(For an algebraic scheme over $k$ having an ample invertible
sheaf is equivalent to being quasi-projective over $k$, see
for example the very general
More on Morphisms, Lemma \ref{more-morphisms-lemma-quasi-projective}.)
By Lemma \ref{lemma-reduced-subgroup-scheme-perfect}
we see that $H$ is a group scheme over $k$.
By Lemma \ref{lemma-reduced-group-scheme-prefect-field-characteristic-p-smooth}
we see that $H$ is smooth over $k$.
This reduces us to the situation discussed in the next
paragraph.

\medskip\noindent
Let $G$ be a quasi-compact irreducible smooth group scheme over an
algebraically closed field $k$. Observe that the local rings of $G$
are regular and hence UFDs
(Varieties, Lemma \ref{varieties-lemma-smooth-regular} and
More on Algebra, Lemma \ref{more-algebra-lemma-regular-local-UFD}).
The complement of a nonempty affine open of $G$
is the support of an effective Cartier divisor $D$.
This follows from Divisors, Lemma
\ref{divisors-lemma-complement-open-affine-effective-cartier-divisor}.
(Observe that $G$ is separated by
Lemma \ref{lemma-group-scheme-over-field-separated}.)
We conclude there exists an effective Cartier divisor $D \subset G$
such that $G \setminus D$ is affine. We will use below that
for any $n \geq 1$ and $g_1, \ldots, g_n \in G(k)$ the complement
$G \setminus \bigcup D g_i$ is affine. Namely, it is the intersection
of the affine opens $G \setminus Dg_i \cong G \setminus D$
in the separated scheme $G$.

\medskip\noindent
We may choose the top row of the diagram
$$
\xymatrix{
G & U \ar[l]_j \ar[r]^\pi & \mathbf{A}^d_k \\
& W \ar[r]^{\pi'} \ar[u] & V \ar[u]
}
$$
such that $U \not = \emptyset$, $j : U \to G$ is an open immersion, and
$\pi$ is \'etale, see
Morphisms, Lemma \ref{morphisms-lemma-smooth-etale-over-affine-space}.
There is a nonempty affine open $V \subset \mathbf{A}^d_k$ such that
with $W = \pi^{-1}(V)$ the morphism $\pi' = \pi|_W : W \to V$ is finite \'etale.
In particular $\pi'$ is finite locally free, say of degree $n$.
Consider the effective Cartier divisor
$$
\mathcal{D} = \{(g, w) \mid m(g, j(w)) \in D\} \subset G \times W
$$
(This is the restriction to $G \times W$ of the pullback of $D \subset G$
under the flat morphism $m : G \times G \to G$.)
Consider the closed subset\footnote{Using the material
in Divisors, Section \ref{divisors-section-norms}
we could take as effective Cartier
divisor $E$ the norm of the effective Cartier divisor $\mathcal{D}$
along the finite locally free morphism $1 \times \pi'$ bypassing
some of the arguments.}
$T = (1 \times \pi')(\mathcal{D}) \subset G \times V$.
Since $\pi'$ is finite locally free, every irreducible component
of $T$ has codimension $1$ in $G \times V$. Since $G \times V$
is smooth over $k$ we conclude these components are effective Cartier
divisors (Divisors, Lemma \ref{divisors-lemma-weil-divisor-is-cartier-UFD}
and lemmas cited above)
and hence $T$ is the support of an effective Cartier divisor
$E$ in $G \times V$. If $v \in V(k)$, then
$(\pi')^{-1}(v) = \{w_1, \ldots, w_n\} \subset W(k)$ and we see that
$$
E_v = \bigcup\nolimits_{i = 1, \ldots, n} D j(w_i)^{-1}
$$
in $G$ set theoretically. In particular we see that $G \setminus E_v$
is affine open (see above).
Moreover, if $g \in G(k)$, then there exists a $v \in V$
such that $g \not \in E_v$. Namely, the set $W'$ of $w \in W$ such that
$g \not \in Dj(w)^{-1}$ is nonempty open and it suffices to pick $v$
such that the fibre of $W' \to V$ over $v$ has $n$ elements.

\medskip\noindent
Consider the invertible sheaf
$\mathcal{M} = \mathcal{O}_{G \times V}(E)$ on $G \times V$.
By Varieties, Lemma \ref{varieties-lemma-rational-equivalence-for-Pic}
the isomorphism class $\mathcal{L}$ of the restriction
$\mathcal{M}_v = \mathcal{O}_G(E_v)$ is independent of $v \in V(k)$.
On the other hand, for every $g \in G(k)$ we can find a $v$
such that $g \not \in E_v$ and such that $G \setminus E_v$
is affine. Thus the canonical section
(Divisors, Definition
\ref{divisors-definition-invertible-sheaf-effective-Cartier-divisor})
of $\mathcal{O}_G(E_v)$
corresponds to a section $s_v$ of $\mathcal{L}$ which does not
vanish at $g$ and such that $G_{s_v}$ is affine.
This means that $\mathcal{L}$ is ample by definition
(Properties, Definition \ref{properties-definition-ample}).
\end{proof}

\bibliographystyle{plain}
\bibliography{my}
\end{document}
